\subsection{مقارنة تعليمات الكتابة وعزل أثر مكوناتها}
\label{sec:prompt-ablation}

تطرح مقارنة تمييز التعليق سؤالًا عن الكتابة: هل يأتي الأسلوب الذي يسهل كشفه من تعليمات الصياغة، أم من التعبير عن موضوع مختار، أم من الجمع بينهما؟ فحصنا هذه المكونات في اختبار تطوير للـ\EN{prompts} قبل اختبار \EN{G5} في تجربة الأرشيف. تغيّر هذه التعليمات كيفية التعبير عن اختيارات التضمين، ولا تقدم ضاغطًا أو مُستقبِلًا مختلفًا.

\subsubsection{عينة تطوير مشتركة وسبعة أنماط للكتابة}
تضم العينة \NUM{\PaPosts} منشورًا من مجموعة المحكّم ذات الـ\NUM{244} منشورًا. اختير \NUM{\PaStratumPosts} منشورًا من كل حالة مشتركة للكشف: كشف الطريقتين، وعدم كشف أي منهما، وكشف طريقتنا وحدها، وكشف \EN{ZLG} وحدها. لذلك فهي عينة تشخيصية موزعة بحسب نتائج معروفة، لا عينة تأكيدية جديدة. وتعيد الأنماط استخدام موضع الرد المختار والموضوع المستهدف، عند وجوده، لكل منشور.

تكوّن الأنماط الأربعة الأولى تصميمًا ذا عاملين، لكل منهما حالتان. يحتفظ \EN{G0} بتعليمات الكتابة التأملية والتعبير عن المتحدث والفعل والشعور، دون موضوع مستهدف. ويمثل \EN{G1} نمط التضمين الأساسي الذي يستخدم تلك التعليمات مع الموضوع المختار. أما \EN{G2}، فيزيل هذا القالب ويطلب فعلًا حواريًا محليًا واحدًا، كسؤال أو تصحيح، دون موضوع مستهدف. ويجمع \EN{G3} التعليمات المحلية مع الموضوع المختار بوصفه تلميحًا لتوجيه الصياغة، لا فكرة يجب إعلانها صراحة. ويساعد النمطان الضابطان \EN{G0} و\EN{G2} على تمييز أثر تعليمات الكتابة من العبء الإضافي للتعبير عن الاختيار المستهدف؛ ولا يمثلان طريقتين لإرسال الرسالة.

تستخدم المقارنة الإضافية العينة نفسها. يضيف \EN{G4} ثلاثة أمثلة بشرية ثابتة للردود إلى \EN{G2}، ويضيفها \EN{G5} إلى \EN{G3}. توجه الأمثلة طول الرد وشكله، مع منع نسخ محتواها. ويطلب نمط إضافي للكتابة الجدلية رأيًا واضحًا مرتبطًا بالمحادثة، مع الإبقاء على الموضوع المستهدف. وهكذا يجمع \EN{G5} تعليمات الرد الحواري الموجز والموضوع المختار والأمثلة البشرية، ولا يمثل خوارزمية تضمين جديدة. اختُبر كل نمط من الأنماط الأربعة الأولى على المنشورات الـ\NUM{\PaPosts} نفسها، بمجموع \NUM{\PaPrimaryAttempts} حالة اختبار. وأضافت الأنماط الثلاثة اللاحقة \NUM{\PaExtensionAttempts} حالة، ليصبح المجموع \NUM{\PaGenerationAttempts}.

\subsubsection{ما الذي اختُبر وما الذي قيّمه المحكّم؟}
يفحص تحقق المُرسِل المعتاد، في الأنماط ذات الموضوع المستهدف، ما إذا كان التعليق المرشح يُقرأ على أنه الموضوع المقصود. ويسجل الحقل المسمى «الاستخراج الدقيق» في هذه التجربة نتيجة تحقق المرشح، إذ لا يجمع برنامجها حمولة كاملة ولا يستخرجها لدى مُستقبِل مستقل. لذلك نسمي النتيجة تحقق الموضوع في الجدول~\ref{tab:prompt-ablation}. ولا ينطبق هذا الفحص على الأنماط الضابطة دون موضوع، ولا يُعد غيابه نسبة استخراج تساوي صفرًا.

تستخدم مهام تمييز التعليق النصوص المقبولة من الأنماط المستهدفة والنصوص المولّدة من الأنماط الضابطة. اختار \EN{GPT-5.6 Luna} بإعداد استدلال مرتفع التعليق المُضاف من مجموعة تضم تسعة ردود بشرية. وتتيح المهام والقرارات المحفوظة مطابقة \NUM{\PaJudgeTasks} حكمًا بنصوص مولّدة غير فارغة. وتختلف أسس عد النصوص المولّدة والمتحقق منها والمعروضة على المحكّم؛ فتبقى الإخفاقات ضمن الحالات المخصصة، بينما يقيّم المحكّم النصوص المؤهلة المتاحة فقط.

\begin{table}[!htbp]
\centering\small
\caption{مقارنة استكشافية للـ\EN{prompts} على \NUM{\PaPosts} منشورًا لكل نمط. يُحسب التوليد وتحقق الموضوع من جميع المنشورات المخصصة، ويُحسب الكشف من النصوص المقيمة فقط، ويكون انخفاضه أفضل في هذه المهمة. تشير الشرطات إلى عدم وجود موضوع مستهدف لفحصه.}
\label{tab:prompt-ablation}
\begin{RTL}
\begin{tabularx}{\columnwidth}{|B|B|B|B|B|}
\hline
النمط & موضوع مستهدف & نص مولّد & تحقق الموضوع & كشف/تقييم \\ \hline
\PaRowsAr
\end{tabularx}
\end{RTL}
\end{table}

\subsubsection{ما الذي توضحه مقارنة التعليمات؟}
لم تؤد إزالة قالب الكتابة التأملية وحدها إلى تغيير عدد حالات الكشف في النمطين الضابطين؛ فكُشف كل من \EN{G0} و\EN{G2} في \NUM{\PaGZeroDetected/\PaGZeroJudged} حالة. وعند وجود موضوع مستهدف، اجتاز \EN{G3} تحقق الموضوع في حالات أكثر من \EN{G1}، بواقع \NUM{\PaGThreeVerified} مقابل \NUM{\PaGOneVerified}، لكن نصوصه المقبولة لم تحقق كشفًا أقل ضمن مجموعة النصوص المقيمة لكل نمط. يميز ذلك بين إنتاج تعليق قابل لقراءة اختياره وجعله يشبه الردود البشرية المحيطة.

اجتاز \EN{G5} تحقق الموضوع في \NUM{\PaGFiveVerified/\PaPosts} حالة مخصصة، وكُشف في \NUM{\PaGFiveDetected/\PaGFiveJudged} نصًا مقيمًا. ولا تعزل نسبة كشفه الأقل من \EN{G3} فائدة الأمثلة البشرية؛ إذ لا توجد نصوص مقبولة تحت النمطين إلا في \NUM{\PaCommonGThreeGFive} منشورًا، وكُشف كل نمط في \NUM{\PaCommonGThreeHits} منها. يرشح التحقق مخرجات مختلفة لكل \EN{prompt}، كما أن إضافة الموضوع المستهدف تضيف شرط قبوله. لذلك تمثل النتائج مقارنات استكشافية للمكونات، ولا تقدّر أثر تعليمة واحدة تقديرًا سببيًا معزولًا أو تثبت الأمان. وحُفظ تكرار منفصل للتوليد مع سجلات تفصيلية، دون دمج نتائجه في الجدول.

يبرر هذا الفحص دراسة \EN{G5} لاحقًا، لكنه لا يجيب عما إذا كانت تعليقاته تستطيع حمل رسالة كاملة متعددة التعليقات. يختبر القسم~\ref{sec:fresh-archive} هذه النتيجة المختلفة؛ إذ لم يستخرج \EN{G5} رسالة كاملة من ست حالات مخصصة، بينما استخرج نمط الكتابة الذي يركز على الموضوع المختار أربع رسائل. ولا تثبت إخفاقات التوليد والتحقق وإخفاق جلب البيانات هناك أن جودة الكتابة أقل. وتوضح الدراستان معًا سبب عدم اعتماد نمط بدا واعدًا في فحص التعليقات بوصفه سياسة الإرسال الأساسية.
